% =============================================================
%  Kapitel 6 – Architektonische Entscheidungen
% =============================================================
\section{Architektonische Entscheidungen}
\label{sec:entscheidungen}

% ------------------------------------------------------------
\subsection{Sprachauswahl: C vs.\ C++}
\label{sec:sprachauswahl}

\begin{center}
\renewcommand{\arraystretch}{1.4}
\begin{tabular}{>{\bfseries}p{3.5cm} p{5.5cm} p{5.5cm}}
  \toprule
  & \textbf{C} & \textbf{C++} \\
  \midrule
  Vorteile
    & Volle Kontrolle, minimaler Overhead, weit verbreitet auf Embedded \\
    & OOP (Klassen, Vererbung), STL, bessere Kapselung, Lesbarkeit \\
  Nachteile
    & Keine Objektorientierung, schlechte Kapselung, fehleranfälliger Code \\
    & Leicht erhöhter Code-Footprint, RTTI/Exceptions ggf.\ deaktivieren \\
  Eignung Arduino Mega
    & Ja
    & Ja (Arduino-Framework basiert auf C++) \\
  Eignung Jetson Nano
    & Ja
    & Ja \\
  \bottomrule
\end{tabular}
\end{center}

\textbf{Entscheidung: C++}

Das gesamte Systemdesign ist objektorientiert modelliert
(Klassen, Interfaces, Enums); C++ erlaubt eine direkte 1:1-Abbildung
dieses UML-Designs in Code ohne strukturelle Kompromisse.
Da das Arduino-Framework selbst auf C++ basiert
(z.\,B.\ \texttt{HardwareSerial}, \texttt{Stream} als C++-Klassen)
und der Jetson Nano Linux mit vollem C++-Support ausführt,
ist C++ auf allen drei Zielplattformen nativ verfügbar.
Zur Minimierung des Speicher- und Laufzeit-Overheads werden
Exceptions und RTTI im Compiler deaktiviert (\texttt{-fno-exceptions
-fno-rtti}), was auf Arduino-Targets Standard ist.
Insgesamt bietet C++ gegenüber reinem C eine deutlich bessere
Kapselung (z.\,B.\ private Methoden in \texttt{EmergencyStop})
bei vernachlässigbarem Mehraufwand.

\textbf{Verworfene Alternativen:}
Reines C würde alle Klassen als Structs mit Funktionszeigern
erfordern – erheblich mehr Boilerplate ohne architektonischen Mehrwert.
Rust wäre sicherheitstechnisch attraktiv, ist aber auf den
verwendeten Arduino-Targets nicht offiziell unterstützt.

% ------------------------------------------------------------
\subsection{Verteilte Architektur}
\label{sec:entsch_verteilung}

\textbf{Entscheidung:} Drei getrennte Recheneinheiten (Jetson, Master, Slave)
statt Mono-Architektur.

\textbf{Rationale:}
\begin{itemize}
  \item \emph{Separation of Concerns:} Wahrnehmung (Jetson) und Sicherheit (Slave)
        sind physisch und softwareseitig getrennt.
  \item \emph{Determinismus:} Der Slave führt eine einfache Regellogik ohne
        OS-Overhead aus.
  \item \emph{Fault Isolation:} Ein Absturz des Jetson deaktiviert den Emergency Stop
        nicht (Fail-Safe: Slave stoppt selbstständig bei ausbleibenden Daten).
\end{itemize}

\textbf{Alternative:} Alles auf dem Jetson Nano (Mono-Architektur).
Wurde verworfen, da dies einen Single Point of Failure schafft und
die Reaktionszeit durch OS-Scheduling unvorhersehbar wird.

% ------------------------------------------------------------
\subsection{Fail-Safe-Prinzip}
\label{sec:entsch_failsafe}

\textbf{Entscheidung:} Im Zweifel immer stoppen.

Konkrete Umsetzungen:
\begin{itemize}
  \item Startup beginnt in \texttt{SoftStop} (kein automatisches Losfahren).
  \item Fehlende Sensordaten $\Rightarrow$ \texttt{SoftStop} + Logging.
  \item \texttt{HardStop} ist von keiner Zustand aus aufhebbar.
\end{itemize}

% ------------------------------------------------------------
\subsection{Klare Verantwortlichkeiten}
\label{sec:entsch_verantwortlichkeiten}

\begin{center}
\renewcommand{\arraystretch}{1.3}
\begin{tabular}{>{\bfseries}p{4cm} p{9.5cm}}
  \toprule
  Komponente & Alleinige Verantwortlichkeit \\
  \midrule
  TTCObserver     & Normal Operation (Gelb/Grün); kein Zugriff auf Emergency Stop. \\
  EmergencyStop   & Sicherheits-Zustände (Rot, Knopf, Taster); kein Zugriff auf TTC. \\
  Logger          & Persistentes Logging; keine Entscheidungslogik. \\
  DriveControlUnit & Motorsteuerung; reagiert nur auf explizite Stop-/Go-Befehle. \\
  \bottomrule
\end{tabular}
\end{center}

% ------------------------------------------------------------
\subsection{CAN als Kommunikationsprotokoll}
\label{sec:entsch_can}

\textbf{Rationale:} CAN ist für Automotive-Systeme optimiert, bietet priorisierungsbasierte
Arbitrierung und eine hardware-gesicherte Fehlerkennung (CRC, ACK). Dies passt besser
zu den Echtzeit-Anforderungen einer Sicherheitsfunktion als eine LAN-basierte Lösung.

\assumption{CAN-Zykluszeit $< 5\,\mathrm{ms}$; konkrete Werte aus Systemtest nachzutragen.}
